\section{Introduction}
\label{sec:introduction}

Kubernetes places a Pod by filtering feasible Nodes and scoring the remaining
alternatives. CPU and memory requests determine placement feasibility, while
resource limits constrain runtime consumption after placement. Once bound, a
running Pod normally remains on its Node until it terminates or is recreated.
The Kubernetes Descheduler complements this behavior by evicting selected Pods
and allowing the kube-scheduler to place their replacements
\citep{k8s_descheduler_docs,larsson_klein_elmroth_2023}.

Placement information alone does not describe the operational cost of an
eviction. A Pod that is an attractive candidate under a consolidation or
load-balancing objective may be processing substantial live demand at the
moment the Descheduler acts. Eviction terminates that instance, after which a
replacement must be scheduled, started, marked Ready, and reattached to
traffic. For a single-replica service, this interval can directly become a
user-visible outage. In this paper, ``single-replica'' specifically means that
the \textit{workload-api} Deployment has one replica Pod behind its Service.

Runtime usage and scheduler feasibility answer different questions. Declared
requests determine where a replacement can fit, whereas current usage
indicates whether disrupting the selected instance is likely to remove active
service capacity. Replacing request-based strategy logic with instantaneous
usage would blur these concerns and make placement decisions sensitive to
short-lived measurements. A more composable approach is to retain the strategy
decision and add a runtime safety gate immediately before eviction.

This paper addresses the following research questions:

\begin{description}
    \item[\textbf{RQ1.}] Does the Actual Usage Evictor prevent
    eviction-induced application failure when a strategy selects a busy
    single-replica API Pod?
    \item[\textbf{RQ2.}] What placement cost does this protection impose under
    consolidation and load-redistribution objectives?
\end{description}

It makes three contributions:

\begin{itemize}
    \item an implementation of the Actual Usage Evictor as a global
    \textit{PreEvictionFilter}, independently composable with Descheduler
    strategies;
    \item a CPU-memory decision model with explicit handling for incomplete
    metrics and missing request baselines; and
    \item a controlled, CPU-path evaluation across the opposing objectives of
    \textit{HighNodeUtilization} and \textit{LowNodeUtilization}, separating
    application protection from strategy-specific placement cost.
\end{itemize}

The evaluation focuses on filter behavior rather than selecting an optimal
threshold. It compares controls using the Default Evictor with treatments that
add the Actual Usage Evictor while holding workload, layout, and host strategy
constant.
